\documentclass[10pt,a4paper]{amsart}
\usepackage[margin=19mm]{geometry}
\usepackage{amsmath,amssymb,mathtools}
\usepackage[scr=rsfs,cal=euler]{mathalfa}
\usepackage{xfrac}
\usepackage[unicode,hidelinks,bookmarks=false]{hyperref}
\newcommand{\ie}{i.e.\ }
\newcommand{\cf}{cf.\ }
\newcommand{\resp}{respectively\ }
\newcommand{\Subsection}{\subsection}
\providecommand{\nobreakdash}{\nobreak\hbox{-}}
\setlength{\emergencystretch}{2em}
\multlinegap0pt
\usepackage{amssymb}
\usepackage{mathtools}
\newtheorem{proposition}{Proposition}
\newtheorem{theorem}{Theorem}
\newcommand{\Imag}{\operatorname{Im}}
\newcommand{\Real}{\operatorname{Re}}
\newcommand{\eps}{\varepsilon}
\title{On the Orthorecursive Expansion of Unity}
\author{Benoit Cloitre}
\date{Source: arXiv:2505.09645v2 (CC BY 4.0)}
\begin{document}
\maketitle

\noindent\emph{Source and licence.} On the Orthorecursive Expansion of Unity. Version arXiv:2505.09645v2 (CC BY 4.0), distributed by arXiv under the Creative Commons Attribution 4.0 International licence, http://creativecommons.org/licenses/by/4.0/, as recorded in the arXiv licence metadata for this exact version and shown on the abstract page https://arxiv.org/abs/2505.09645v2. Full licence notice: \texttt{licenses/arxiv-2505.09645-CC-BY-4.0.txt}.\par\smallskip

\noindent\textit{Editorial scope.} Counted unit: the article's theorem thm:zerofree (source0.tex lines 391--393). The article's complete original proof of it is delivered with it (source0.tex lines 395--431, one \texttt{proof} environment of 37 lines). No mathematics is written, repaired or re-derived: the statement and the proof are the article's own lines, in the article's own notation, and no retained line is substituted or re-flowed.  Retained uncounted as the source-local context the proof needs: lines 339--341; lines 348--350; lines 351--353; lines 376--378.  Nothing was omitted from any retained span, so the package carries no omission record.  No retained line contains a citation, so the wrapper carries no bibliography and no bibliography entry.  No retained line contains a figure environment, an image-inclusion command or drawing code, so no image is delivered and the package declares no asset dependency.  Editorial formatting only, in addition: an AMS \texttt{amsart} preamble that restates the article's own declarations and macros used by the retained lines, namely the environments \texttt{proposition}, \texttt{theorem} on separate counters, together with the standard packages those lines need.  The retained environments print in this document as Proposition 1, Theorem 1 in order of appearance, whereas the article prints the same items with the numbering of its own full document.  The complete native source of the article as distributed in the arXiv e-print for this exact version is bundled unchanged as \texttt{sources/178130/source0.tex}.
\begin{equation}\label{eq:Dinfty_asymp}
D_\infty(z) = \frac{1}{2z} - \frac{1}{4z^2} + O\!\left(\frac{1}{|z|^3}\right) \qquad (|\Imag(z)| \to \infty)
\end{equation}

\begin{equation}\label{eq:real_part}
\Real\, D_\infty(z) = \sum_{j=0}^{\infty} \frac{(-1)^j(x-j)}{(x-j)^2 + y^2},
\end{equation}

\begin{equation}\label{eq:imag_part}
\Imag\, D_\infty(z) = -y \sum_{j=0}^{\infty} \frac{(-1)^j}{(x-j)^2 + y^2}.
\end{equation}

\begin{proposition}\label{prop:zerofree_left}
For $z = x + iy$ with $x \leq 0$ and $y \neq 0$, one has $\Imag\, D_\infty(z) \neq 0$. In particular, $D_\infty(z) \neq 0$ for all $z$ with $\Real(z) \leq 0$.
\end{proposition}

\begin{theorem}[Zero-free strip]\label{thm:zerofree}
The function $D_\infty$ has no zeros in the closed strip $\{z \in \mathbb{C} : 0 < \Real(z) \leq 1,\; \Imag(z) > 0\}$, and therefore $g^*(z) \neq 0$ for $\Real(z) \leq 1$.
\end{theorem}

\begin{proof}
We apply the argument principle to $D_\infty$ on the positively oriented contour $\Gamma_T$ bounding the region $\{0 < \Real(z) \leq 1,\; 0 < \Imag(z) < T\}$, with small indentations of radius $\eps$ around the poles at $z = 0$ and $z = 1$. The contour consists of six segments traversed in order. We show that the total variation of $\arg D_\infty(z)$ around $\Gamma_T$ is zero for $T$ large enough.

Along the bottom segment ($z = x$ for $0 < x < 1$), the absence-of-real-zeros analysis with $p = 0$ gives $D_\infty(x) > 0$, so the argument does not change.

Near $z = 0$, one has $D_\infty(z) = 1/z + h_0(z)$ where $h_0$ is bounded. On the quarter-circle $z = \eps e^{i\theta}$ with $\theta$ from $\pi/2$ to $0$, the argument of $D_\infty(z)$ is approximately $-\theta$. The variation as $\eps \to 0$ is $+\pi/2$. Near $z = 1$, the residue is $-1$, so $D_\infty(z) \sim -1/(z-1)$. On the quarter-circle $z = 1 + \eps e^{i\theta}$ with $\theta$ from $\pi$ to $\pi/2$, the argument is approximately $\pi - \theta$, giving a variation of $+\pi/2$.

Along the left wall ($z = iy$ with $y > 0$), the same alternating-series argument as in Proposition~\ref{prop:zerofree_left} shows that $\Imag\, D_\infty(iy) < 0$ for all $y > 0$. The image stays in the lower half-plane, so the argument does not wind.

Along the ceiling ($z = x + iT$ with $0 \leq x \leq 1$), the asymptotic~\eqref{eq:Dinfty_asymp} gives $D_\infty(x+iT) = 1/(2(x+iT)) + O(T^{-2})$. For $T$ large enough, the image is a small perturbation of a curve in the fourth quadrant. The argument variation is $O(1/T) \to 0$.

It remains to analyze the right wall ($z = 1 + iy$ with $y > 0$), which is the most delicate segment. We treat two ranges of $y$ separately.

For $0 < y \leq 2$, we show $\Real\, D_\infty(1+iy) > 0$. By~\eqref{eq:real_part} with $x = 1$, the $j=0$ term is $1/(1+y^2)$, the $j=1$ term vanishes (numerator $1-1=0$), and for $j \geq 2$ we substitute $m = j-1$ to get
$$
\Real\, D_\infty(1+iy) = \sum_{m=2}^{\infty} \frac{(-1)^m m}{m^2 + y^2}.
$$
The cancellation of the $j=0$ term with the $m=1$ term has already been incorporated. The function $t \mapsto t/(t^2 + y^2)$ is strictly decreasing for $t \geq y$. Since $m \geq 2 \geq y$ (as $y \leq 2$), the terms strictly decrease in magnitude. By Leibniz's criterion, the sum exceeds its first two terms:
$$
\Real\, D_\infty(1+iy) > \frac{2}{4+y^2} - \frac{3}{9+y^2} = \frac{6-y^2}{(4+y^2)(9+y^2)}.
$$
For $0 < y \leq 2$, one has $y^2 \leq 4 < 6$, so this is strictly positive.

For $y \geq 2$, we show $\Imag\, D_\infty(1+iy) < 0$. By~\eqref{eq:imag_part} with $x = 1$, the terms for $j = 0,1,2$ yield
$$
\frac{2}{1+y^2} - \frac{1}{y^2} = \frac{y^2-1}{y^2(1+y^2)}.
$$
The tail starting at $j = 3$ is an alternating series with strictly decreasing terms, bounded below by its first term $-1/(4+y^2)$. Adding and simplifying,
$$
\frac{\Imag\, D_\infty(1+iy)}{-y} \geq \frac{y^2-1}{y^2(1+y^2)} - \frac{1}{4+y^2} = \frac{2y^2-4}{y^2(1+y^2)(4+y^2)} > 0
$$
for $y > \sqrt{2}$, hence for $y \geq 2$.

Combining both ranges, $D_\infty(1+iy) \neq 0$ for all $y > 0$, and the image of the right wall avoids the negative real axis. As $y \to 0^+$, $D_\infty(1+iy) \sim -1/(iy) \to +i\infty$, so the argument starts at $\pi/2$. As $y$ increases, the image passes through the right half-plane (for $y \leq 2$), then into the lower half-plane (for $y \geq 2$), and ends near $-i/(2T)$ in the fourth quadrant (argument approximately $-\pi/2$). The total variation along the right wall is $-\pi$.

Summing over all segments, $\Delta\arg = 0 + \pi/2 + (-\pi) + 0 + 0 + \pi/2 = 0$. There are no zeros of $D_\infty$ in the closed strip $0 < \Real(z) \leq 1$, $\Imag(z) > 0$.
\end{proof}

\end{document}
